% =============================================================
%  hydrothermal_goethite_body.tex   --   REBUILD v2
%
%  MERGED PAPER (Chapters 1 + 2), goethite-led, targeted at
%  Geochimica et Cosmochimica Acta.
%
%  Chapters 1 and 2 remain the extended thesis record and are NOT
%  edited by this file.
%
%  WHAT CHANGED FROM v1 (per BJT, this rebuild)
%  --------------------------------------------
%  The paper is now built around a fixed figure sequence, and the
%  empty- and sulfur-ampoule calibration experiments are GONE: not
%  named, not characterized, not tabulated. Only one sentence in the
%  Discussion notes that water-bearing ampoules show anomalously high
%  rate constants, as motivation for measuring the rate constant in
%  situ.
%
%  This is a structural improvement, not just a cut. In v1 the
%  conservative-bound argument leaned on a separate water calibration
%  ampoule. Here it rests on the goethite ampoule's OWN measured rate
%  constant exceeding the anhydrous literature value (Fig. 3), which
%  is self-contained, needs no companion experiment, and is a stronger
%  argument besides.
%
%  ARC (BJT's four moves):
%    1. The ampoule stabilizes reactive phases during step heating,
%       which motivates investigating ampoule systematics.  [Fig. 1]
%    2. Ampoule systematics are tractable: encapsulated Durango
%       apatite recovers its published Arrhenius array.       [Fig. 2]
%    3. The permeation theory, demonstrated by measuring the rate
%       constant of the goethite ampoule itself. This is the bridge
%       from Chapter 1 into Chapter 2.                 [Figs. 3, 4; Tables 1, 2]
%    4. Goethite results, discussion, science.          [Figs. 5-9]
%
%  FIGURE / TABLE SEQUENCE (fixed by BJT):
%    Fig. 1  stabilization grid, 3 rows (pyrite / arsenopyrite / goethite)
%    Fig. 2  encapsulated_DApErr.pdf
%    Fig. 3  NaiveIsothermalFits_Grid.pdf   (low-T sequence fits)
%    Fig. 4  NaiveLambda_Arrhenius.pdf
%    Table 1 Winsor 3He data      \begin{landscape}
\begin{footnotesize}
\begin{longtable}{r r r r r r r r r r}
\caption{$^3$He step-heat dataset for HT-Win-06-01b under the ``treatment as uncertainty'' method: measured and cumulative $^3$He, the upper-bound residual ampoule content ($n_{t,\mathrm{UB}}$), the lower- and upper-bound true step production ($\omega_{\mathrm{LB}}$, $\omega_{\mathrm{UB}}$), and the corresponding lower-bound, central, and upper-bound apparent diffusivities. ``--'' denotes early steps with $\lambda t < 4$, for which bounds could not be determined (see \S\ref{sec:rate-constant}). Six rows have derived values marked with ``*'': these steps have measured $^3$He $<0.1$ pccSTP and are excluded as low-signal from all plots (Fig.~\ref{fig:ratio-spectra} and others derived from this dataset).}\label{tab:winsor-3he}\\
\toprule
Temp & Time & $^3$He meas. & Cum. $^3$He & $n_{t,\mathrm{UB}}$ & $\omega_{\mathrm{LB}}$ & $\omega_{\mathrm{UB}}$ & $\ln(D/a^2)_{\mathrm{LB}}$ & $\ln(D/a^2)_{m}$ & $\ln(D/a^2)_{\mathrm{UB}}$ \\
(°C) & (min) & (pccSTP) & (pccSTP) & (pccSTP) & (pccSTP) & (pccSTP) &  &  &  \\
\midrule
\endfirsthead
\multicolumn{10}{l}{\textit{(continued)}} \\
\toprule
Temp & Time & $^3$He meas. & Cum. $^3$He & $n_{t,\mathrm{UB}}$ & $\omega_{\mathrm{LB}}$ & $\omega_{\mathrm{UB}}$ & $\ln(D/a^2)_{\mathrm{LB}}$ & $\ln(D/a^2)_{m}$ & $\ln(D/a^2)_{\mathrm{UB}}$ \\
(°C) & (min) & (pccSTP) & (pccSTP) & (pccSTP) & (pccSTP) & (pccSTP) &  &  &  \\
\midrule
\endhead
\midrule
\multicolumn{10}{r}{\textit{continued on next page}} \\
\endfoot
\bottomrule
\endlastfoot
170 & 300 & 0.9630 & 0.0000 & 1.1660 & -- & -- & -- & -- & -- \\
180 & 463 & 0.5442 & 0.9630 & 0.2413 & -- & -- & -- & -- & -- \\
190 & 550 & 0.5292 & 1.6241 & 0.1538 & -- & -- & -- & -- & -- \\
200 & 500 & 0.5190 & 2.2112 & 0.1430 & 0.3651 & 0.6619 & -22.158 & -21.775 & -21.503 \\
210 & 475 & 0.5770 & 2.7646 & 0.1440 & 0.4341 & 0.7210 & -21.715 & -21.406 & -21.160 \\
220 & 450 & 0.6333 & 3.3799 & 0.1447 & 0.4894 & 0.7781 & -21.345 & -21.068 & -20.843 \\
230 & 425 & 0.7899 & 4.0551 & 0.1685 & 0.6452 & 0.9584 & -20.823 & -20.604 & -20.392 \\
225 & 908 & 0.7513 & 4.8885 & 0.0716 & 0.5827 & 0.8229 & -21.516 & -21.246 & -21.148 \\
215 & 450 & 0.1415 & 5.6842 & 0.0349 & 0.0699 & 0.1764 & -22.836 & -22.124 & -21.900 \\
205 & 520 & 0.0763 & 5.8529 & 0.0184$^*$ & 0.0413$^*$ & 0.0947$^*$ & -23.479$^*$ & -22.863$^*$ & -22.646$^*$ \\
195 & 1106 & 0.0670 & 5.9491 & 0.0080$^*$ & 0.0487$^*$ & 0.0750$^*$ & -24.053$^*$ & -23.731$^*$ & -23.618$^*$ \\
185 & 740 & 0.0227 & 6.0246 & 0.0049$^*$ & 0.0147$^*$ & 0.0276$^*$ & -24.839$^*$ & -24.404$^*$ & -24.207$^*$ \\
175 & 920 & 0.0180 & 6.0610 & 0.0036$^*$ & 0.0131$^*$ & 0.0216$^*$ & -25.168$^*$ & -24.847$^*$ & -24.665$^*$ \\
200 & 260 & 0.0224 & 6.0814 & 0.0163$^*$ & 0.0189$^*$ & 0.0388$^*$ & -23.534$^*$ & -23.360$^*$ & -22.812$^*$ \\
210 & 520 & 0.0812 & 6.1081 & 0.0181$^*$ & 0.0649$^*$ & 0.0993$^*$ & -22.984$^*$ & -22.758$^*$ & -22.554$^*$ \\
220 & 1167 & 0.5766 & 6.1914 & 0.0447 & 0.5585 & 0.6213 & -21.587 & -21.553 & -21.475 \\
230 & 370 & 0.3577 & 6.7796 & 0.0903 & 0.3130 & 0.4480 & -20.948 & -20.811 & -20.579 \\
240 & 320 & 0.6767 & 7.1563 & 0.1776 & 0.5864 & 0.8543 & -20.062 & -19.913 & -19.667 \\
245 & 250 & 0.5549 & 7.8582 & 0.1888 & 0.3773 & 0.7437 & -20.176 & -19.778 & -19.474 \\
250 & 210 & 0.6934 & 8.4743 & 0.2786 & 0.5047 & 0.9721 & -19.626 & -19.296 & -18.942 \\
245 & 270 & 0.4659 & 9.2086 & 0.1413 & 0.1872 & 0.6071 & -20.801 & -19.874 & -19.601 \\
240 & 320 & 0.2857 & 9.7206 & 0.0749 & 0.1445 & 0.3606 & -21.176 & -20.487 & -20.250 \\
235 & 370 & 0.2199 & 10.0357 & 0.0517 & 0.1450 & 0.2716 & -21.284 & -20.864 & -20.650 \\
230 & 420 & 0.1400 & 10.2736 & 0.0303 & 0.0883 & 0.1703 & -21.885 & -21.421 & -21.224 \\
225 & 652 & 0.1494 & 10.4279 & 0.0206 & 0.1191 & 0.1700 & -22.009 & -21.781 & -21.650 \\
230 & 500 & 0.1550 & 10.5908 & 0.0271 & 0.1344 & 0.1820 & -21.605 & -21.462 & -21.299 \\
235 & 450 & 0.1951 & 10.7582 & 0.0361 & 0.1681 & 0.2313 & -21.258 & -21.107 & -20.936 \\
240 & 400 & 0.2425 & 10.9638 & 0.0483 & 0.2063 & 0.2908 & -20.913 & -20.750 & -20.566 \\
245 & 350 & 0.3150 & 11.2328 & 0.0690 & 0.2667 & 0.3840 & -20.495 & -20.326 & -20.125 \\
250 & 300 & 0.3716 & 11.5689 & 0.0939 & 0.3026 & 0.4656 & -20.181 & -19.973 & -19.743 \\
250 & 300 & 0.3640 & 11.9627 & 0.0899 & 0.2700 & 0.4539 & -20.261 & -19.959 & -19.734 \\
250 & 300 & 0.3452 & 12.3490 & 0.0852 & 0.2553 & 0.4304 & -20.285 & -19.979 & -19.755 \\
250 & 10250 & 8.4829 & 12.7167 & 0.0494 & 8.3977 & 8.5324 & -19.989 & -19.976 & -19.969 \\
240 & 785 & 0.2261 & 21.2187 & 0.0209 & 0.1767 & 0.2470 & -21.026 & -20.778 & -20.689 \\
\midrule
\textbf{Total} &  &  & 146.1483 &  &  &  &  &  &  \\
\end{longtable}
\end{footnotesize}
\end{landscape}

%    Table 2 Winsor 4He data      \begin{landscape}
\begin{footnotesize}
\begin{longtable}{r r r r r r r r r r r r r}
\caption{$^4$He step-heat dataset and absolute $^4$He/$^3$He ratios for HT-Win-06-01b under the ``treatment as uncertainty'' method, structured as in Table~\ref{tab:winsor-3he}. ``--'' denotes steps for which bounds could not be determined. The Total row reports the bulk (total-yield) $^4$He/$^3$He ratio, i.e. the normalizing value used to bulk-normalize these ratios for the plot in Fig.~\ref{fig:ratio-spectra}. Six rows have derived values marked with ``*'': these steps have measured $^3$He $<0.1$ pccSTP and are excluded as low-signal from all plots, as in Table~\ref{tab:winsor-3he}.}\label{tab:winsor-4he}\\
\toprule
Temp & Time & $^4$He meas. & Cum. $^4$He & $n_{t,\mathrm{UB}}$ & $\omega_{\mathrm{LB}}$ & $\omega_{\mathrm{UB}}$ & $\ln(D/a^2)_{\mathrm{LB}}$ & $\ln(D/a^2)_{m}$ & $\ln(D/a^2)_{\mathrm{UB}}$ & $(^4\mathrm{He}/^3\mathrm{He})_{\mathrm{LB}}$ & $(^4\mathrm{He}/^3\mathrm{He})_{m}$ & $(^4\mathrm{He}/^3\mathrm{He})_{\mathrm{UB}}$ \\
(°C) & (min) & (nccSTP) & (nccSTP) & (nccSTP) & (nccSTP) & (nccSTP) &  &  &  &  &  &  \\
\midrule
\endfirsthead
\multicolumn{13}{l}{\textit{(continued)}} \\
\toprule
Temp & Time & $^4$He meas. & Cum. $^4$He & $n_{t,\mathrm{UB}}$ & $\omega_{\mathrm{LB}}$ & $\omega_{\mathrm{UB}}$ & $\ln(D/a^2)_{\mathrm{LB}}$ & $\ln(D/a^2)_{m}$ & $\ln(D/a^2)_{\mathrm{UB}}$ & $(^4\mathrm{He}/^3\mathrm{He})_{\mathrm{LB}}$ & $(^4\mathrm{He}/^3\mathrm{He})_{m}$ & $(^4\mathrm{He}/^3\mathrm{He})_{\mathrm{UB}}$ \\
(°C) & (min) & (nccSTP) & (nccSTP) & (nccSTP) & (nccSTP) & (nccSTP) &  &  &  &  &  &  \\
\midrule
\endhead
\midrule
\multicolumn{13}{r}{\textit{continued on next page}} \\
\endfoot
\bottomrule
\endlastfoot
170 & 300 & 0.0192 & 0.0000 & 0.0280 & -- & -- & -- & -- & -- & -- & -- & -- \\
180 & 463 & 0.0305 & 0.0192 & 0.0152 & -- & -- & -- & -- & -- & -- & -- & -- \\
190 & 550 & 0.0558 & 0.0660 & 0.0179 & -- & -- & -- & -- & -- & -- & -- & -- \\
200 & 500 & 0.1024 & 0.1375 & 0.0313 & 0.0845 & 0.1337 & -26.207 & -25.965 & -25.619 & -- & 197.3 & -- \\
210 & 475 & 0.1460 & 0.2500 & 0.0404 & 0.1147 & 0.1864 & -25.314 & -25.022 & -24.717 & 248.2 & 253.0 & 268.8 \\
220 & 450 & 0.2091 & 0.4055 & 0.0530 & 0.1686 & 0.2620 & -24.407 & -24.152 & -23.876 & 326.3 & 330.1 & 349.3 \\
230 & 425 & 0.3243 & 0.6319 & 0.0760 & 0.2714 & 0.4004 & -23.425 & -23.213 & -22.956 & 403.7 & 410.6 & 432.3 \\
225 & 908 & 0.3483 & 0.9799 & 0.0363 & 0.2722 & 0.3845 & -23.807 & -23.527 & -23.413 & 442.6 & 463.6 & 474.4 \\
215 & 450 & 0.0716 & 1.3487 & 0.0196 & 0.0353 & 0.0912 & -24.945 & -24.225 & -23.976 & 429.1 & 505.9 & 537.1 \\
205 & 520 & 0.0385 & 1.4359 & 0.0103$^*$ & 0.0189$^*$ & 0.0487$^*$ & -25.661$^*$ & -24.942$^*$ & -24.701$^*$ & 392.4$^*$ & 504.2$^*$ & 534.6$^*$ \\
195 & 1106 & 0.0336 & 1.4912 & 0.0043$^*$ & 0.0233$^*$ & 0.0379$^*$ & -26.165$^*$ & -25.796$^*$ & -25.675$^*$ & 446.7$^*$ & 500.9$^*$ & 515.3$^*$ \\
185 & 740 & 0.0079 & 1.5467 & 0.0019$^*$ & 0.0037$^*$ & 0.0098$^*$ & -27.586$^*$ & -26.808$^*$ & -26.592$^*$ & 224.7$^*$ & 350.1$^*$ & 369.1$^*$ \\
175 & 920 & 0.0152 & 1.6020 & 0.0033$^*$ & 0.0133$^*$ & 0.0185$^*$ & -26.475$^*$ & -26.340$^*$ & -26.141$^*$ & 843.2$^*$ & 843.2$^*$ & 982.0$^*$ \\
200 & 260 & 0.0107 & 1.6351 & 0.0091$^*$ & 0.0073$^*$ & 0.0198$^*$ & -25.787$^*$ & -25.411$^*$ & -24.791$^*$ & 375.6$^*$ & 475.5$^*$ & 554.7$^*$ \\
210 & 520 & 0.2166 & 1.6547 & 0.0535$^*$ & 0.2075$^*$ & 0.2701$^*$ & -23.067$^*$ & -23.021$^*$ & -22.786$^*$ & 2667.1$^*$ & 2667.1$^*$ & 3103.8$^*$ \\
220 & 1167 & 0.3187 & 1.8714 & 0.0269 & 0.2652 & 0.3456 & -23.499 & -23.302 & -23.214 & 471.7 & 552.7 & 563.2 \\
230 & 370 & 0.2113 & 2.6766 & 0.0590 & 0.1844 & 0.2703 & -22.391 & -22.249 & -21.993 & 575.1 & 590.8 & 628.0 \\
240 & 320 & 0.4286 & 3.1616 & 0.1243 & 0.3697 & 0.5529 & -21.360 & -21.203 & -20.931 & 612.1 & 633.4 & 674.8 \\
245 & 250 & 0.3988 & 3.6430 & 0.1521 & 0.2746 & 0.5510 & -21.289 & -20.899 & -20.556 & 665.9 & 718.7 & 778.4 \\
250 & 210 & 0.4440 & 4.2359 & 0.2020 & 0.2918 & 0.6460 & -20.906 & -20.469 & -20.072 & 541.7 & 640.2 & 702.8 \\
245 & 270 & 0.3123 & 4.7263 & 0.1053 & 0.1103 & 0.4176 & -22.043 & -20.981 & -20.680 & 469.1 & 670.4 & 721.4 \\
240 & 320 & 0.1927 & 5.0846 & 0.0559 & 0.0874 & 0.2485 & -22.376 & -21.575 & -21.315 & 505.5 & 674.3 & 718.4 \\
235 & 370 & 0.1517 & 5.3031 & 0.0393 & 0.0958 & 0.1910 & -22.386 & -21.922 & -21.688 & 599.8 & 689.6 & 729.9 \\
230 & 420 & 0.0991 & 5.4714 & 0.0236 & 0.0598 & 0.1227 & -22.956 & -22.448 & -22.232 & 608.3 & 707.9 & 745.9 \\
225 & 652 & 0.0995 & 5.5835 & 0.0151 & 0.0760 & 0.1146 & -23.136 & -22.863 & -22.721 & 608.2 & 666.5 & 689.2 \\
230 & 500 & 0.1051 & 5.6969 & 0.0202 & 0.0900 & 0.1253 & -22.679 & -22.523 & -22.345 & 652.1 & 678.1 & 707.5 \\
235 & 450 & 0.1369 & 5.8251 & 0.0279 & 0.1167 & 0.1648 & -22.290 & -22.128 & -21.940 & 674.6 & 701.7 & 734.0 \\
240 & 400 & 0.1781 & 5.9762 & 0.0390 & 0.1502 & 0.2171 & -21.892 & -21.719 & -21.518 & 705.1 & 734.6 & 770.9 \\
245 & 350 & 0.2305 & 6.1807 & 0.0556 & 0.1915 & 0.2861 & -21.479 & -21.290 & -21.070 & 694.4 & 731.9 & 771.7 \\
250 & 300 & 0.2777 & 6.4385 & 0.0758 & 0.2221 & 0.3535 & -21.134 & -20.906 & -20.659 & 703.3 & 747.3 & 793.3 \\
250 & 300 & 0.2654 & 6.7374 & 0.0725 & 0.1896 & 0.3379 & -21.250 & -20.908 & -20.661 & 653.7 & 729.2 & 774.1 \\
250 & 300 & 0.2489 & 7.0231 & 0.0680 & 0.1764 & 0.3169 & -21.282 & -20.932 & -20.647 & 646.8 & 721.0 & 765.4 \\
250 & 10250 & 6.8042 & 7.2901 & 0.0430 & 6.7362 & 6.8472 & -20.704 & -20.691 & -20.682 & 800.6 & 802.1 & 803.3 \\
240 & 785 & 0.1853 & 14.1156 & 0.0187 & 0.1423 & 0.2040 & -21.688 & -21.422 & -21.325 & 770.5 & 819.4 & 838.0 \\
\midrule
\textbf{Total} &  &  & 145.8145 &  &  &  &  &  &  &  & 997.7 &  \\
\end{longtable}
\end{footnotesize}
\end{landscape}

%    Fig. 5  zoomout2.jpg
%    Fig. 6  winsor_goethite_Diffusivities_Arrhenius.pdf
%    Fig. 7  winsor_goethite_Diffusivities_Arrhenius_Zoomin.pdf
%    Fig. 8  DeltasWinsorGTzoomin_VacuumHydrothermal.pdf
%    Fig. 9  Comparative_Ratios_Dtplot_WinsorGT_VacuumHDMKAF.pdf
%    SI      Durango step heat table  -> hydrothermal_goethite_SI_body.tex
%
%  NOTE ON FIG. 3 / FIG. 4 ORDER: these were swapped from BJT's
%  original sequence, which had the Arrhenius summary ahead of the
%  sequence fits that produce it. They now run in causal order
%  (fits -> rate constants -> Arrhenius), which matches the order the
%  prose in \S\ref{sec:results-rate-constant} already followed.
%
%  Fig. 4 currently plots measured permeabilities K. When BJT recasts
%  it in terms of the rate constant lambda, the following need to
%  change together:
%    - Fig. 4 caption: drop "converted from the rate constants fitted
%      in Fig. 3 via K = lambda/GRT"; the conversion disappears.
%    - Fig. 3 caption, last sentence: same, "converted to
%      permeabilities ... plotted against temperature" becomes simply
%      "plotted against temperature in Fig. 4".
%    - \S\ref{sec:results-rate-constant} para 2 opener: "Converting
%      these rate constants to permeabilities through K = lambda/GRT"
%      becomes "Every fitted rate constant ...", and "regression of
%      ln K" becomes "regression of ln lambda".
%    - The activation energy is NOT invariant under the recast, so
%      \ChTwoPermEa must be regenerated, not reused. Since
%      lambda = K G R T,  ln lambda = ln K + ln(GR) + ln T, and
%      d(ln T)/d(1/T) = -T, the fitted slope shifts by exactly RT:
%          Ea(lambda) = Ea(K) + RT
%      At the ~373 K midpoint of the 90-110 C fitting range that is
%      about +740 cal/mol, i.e. ~7% of the current 10192 cal/mol.
%      That is inside the current +/-1882 (2 s.e.) uncertainty, but
%      it is a real and predictable offset: use it to check that the
%      regenerated regression is right rather than quoting the old
%      number.
%
%  PORTABILITY
%    Class-agnostic body: no \documentclass, no \begin{document}, no
%    abstract wrapper. \input{} from a journal or thesis wrapper.
%    \graphicspath must include:
%      Chapter1/Figures/generated/   Chapter1/Figures/source/
%      Chapter2/Figures/generated/   Chapter2/Figures/source/
%    Requires: booktabs, longtable, graphicx, amsmath, lscape (or
%    pdflscape) for the landscape data tables.
%    \input Chapter2/Writing/chapter2_derived_values.tex for \ChTwo* macros.
% =============================================================

\providecommand{\tophline}{\toprule}
\providecommand{\middlehline}{\midrule}
\providecommand{\bottomhline}{\bottomrule}

% =============================================================
\section{Introduction}\label{sec:introduction}
% =============================================================

Goethite (FeOOH) crystallization is datable only by U-Th/He geochronology \citep{shustervasconcelos2005,farley2024}. Goethite is a powerful chronometer for dating surface processes such weathering and ferruginization directly: laterite and duricrust formation and recrystallization \citep{monteiro2014,allard2018}, and the ferruginization of organic matter and paleosols \citep{heim2006,yappshuster2011,vasconcelos2013,danisik2013,hofmann2017}. 

Interpreting a goethite U-Th/He date requires knowing the conditions under which the mineral retains radiogenic helium. Step heats of proton-irradiated samples measure the degree of partial helium loss by the \textsuperscript{4}He/\textsuperscript{3}He method \citep{shustervasconcelos2005}, and reported losses are always non-zero but modest, ranging from 2 to 30\% \citep{shustervasconcelos2005,heim2006,hofmann2017,deng2017}, but the effects of temperature, crystallite size, and chemical composition have not been determined experimentally and remain unresolved because goethite (FeOOH) decomposes to Hematite (Fe\textsubscript{2}O\textsubscript{2}) in the step heat experiments used to characterize it \citep{farley2024}. \citet{farley2024} measured evolved water by capacitance manometry during vacuum step heats and found helium and water release to be tightly correlated even below 200\textdegree C, indicating that decomposition, rather than helium migration through an unchanging crystal, governs the release. Vacuum step-heat results for goethite therefore need not describe helium behavior in nature.

Goethite dehydration is suppressed at high water pressure, so a step heat conducted under hydrothermal conditions may interrogate helium loss rather than decomposition kinetics. 
A particularly simple and agile approach to do this is to seal the sample inside a pure silica glass ampoule. Silica glass is the most helium-permeable common glass \citep{norton1953,altemose1961} because it lacks the network-modifying cations that occupy interstitial space in other glasses. Helium permeation in silica glass is sufficiently rapid, so helium evolved from a sample inside the ampoule escapes to the vacuum system and can be measured while other gaseous species, such as water, sulfur, or oxygen is retained. Earlier work has exploited the permeability and selectivity of silica glass for U-Th/He geochronology \citep{yakubovich2020,yakubovich2021,yakubovich2024} and for monotonic ramped-heating experiments \citep{yakubovich2019}, but not to impose hydrothermal conditions during step heating. 

Silica glass encapsulation offers a unique opportunity to preserve reactive phases such as goethite that would otherwise decompose during helium step heat experiments, motivating the characterization of a step heat method utilizing this approach. Early applications reveal one cost: Permeation through the ampoule wall is rapid but not instantaneous, so helium released during a heating step is measured partly during that step and partly during later ones. We refer to this offset as permeation lag, and it must be quantified before any encapsulated step heat can be faithfully interpreted.

This paper establishes that the ampoule contribution to the measured helium in step heat experiments is tractable and that this method can address questions that vacuum experiments cannot. We show that encapsulation stabilizes reactive phases during heating and that we can successfully deconvolve the ampoule's permation lag effect from helium step heat measurements of Durango Apatite to recover its well-documented helium diffusion systematics \citep{wolf1998, farley2000}. We show that the ampoule's permeation rate constant can be measured during the experiment itself, which is necessary to apply the method to goethite because water-bearing ampoules often permeate heium anomalously fast. Finally, we apply the method to proton-irradiated goethite Win-06-01b under water saturation at 170--250\textdegree C, the first hydrothermal helium step heat reported for goethite.

% =============================================================
\section{Materials and Methods}\label{sec:methods}
% =============================================================

\subsection{Ampoule preparation and sealing}\label{sec:ampoule-prep}

Segments of pure silica glass tube approximately 10\,cm long, 4\,mm outer diameter and 3\,mm inner diameter, were sonicated in sequential baths of acetone, ethanol, and MilliQ water, then dried overnight at 100\textdegree C. The tubes were hand-pulled into pipettes with a jewelry welding torch and the thin ends were fused closed. The sample and buffering phase were then loaded into the pipette tip. Inner and outer diameters were measured on a calibrated stereo microscope. These give the geometric factor \emph{G} of Eq.~\eqref{eq:geomfactor}: 2393.5\,cm\textsuperscript{-2} for the Durango apatite ampoule (850\,\textmu m inner, 1350\,\textmu m outer diameter) and, for the goethite ampoule, the value used throughout \S\ref{sec:results-rate-constant}.

It is paramount to avoid heating the sample during the encapulation process to preserve the integrity of the step heat experimentl however, this challenging because silica glass softens well above 1600\textdegree C. To protect against heating during the sealing process, the sample-end of the ampoule was held by a water-soaked Kimwipe while the top was fused. Liquid water saturated the wipe afterward, so the adjacent glass containing the sample never exceeded 100\textdegree C. For goethite, the contents were frozen and held in liquid nitogen instead, serving the same purpose while additionally preventing the loss of the water charge.

Before sealing, each ampoule was pumped to below 0.1\,torr through its open top. The purpose of evacuation is mechanical: under an internal-external pressure contrast, softening glass deforms inward and self-seals, which gives a better seal with far less heat input. Unfortunately, this process may not usefully reduce the ampoule's helium inventory, because room-temperature permeation is fast enough that atmospheric helium re-enters a sealed ampoule within hours to days. The residual atmospheric helium is instead measured directly at the start of each experiment (\S\ref{sec:permeation-steps}).

Ampoules containing water withstood the vapor pressure of liquid-water saturation up to about 86\,bar at 300\textdegree C \citep{haar1984} but ruptured above 315\textdegree C. We recommend limiting water-bearing ampoules to below 300\textdegree C.

\subsection{Encapsulated Durango apatite}\label{sec:methods-apatite}

Durango apatite crushate sieved to 300\,\textmu m (radii $\approx$ 150\,\textmu m), 3.18\,mg, was encapsulated and step heated over 100--400\textdegree C. Durango apatite is a standard for helium diffusion and U-Th/He thermochronometry \citep{zeitler1987,wolf1996,farley2000,shuster2005,shuster2006,flowers2007}, so its Arrhenius behavior is known independently. The experiment then is a direct test of the validity of the permeation-lag correction.

Each high-temperature step was preceded and followed by a sequence of two or three low-temperature permeation steps (\S\ref{sec:permeation-steps}). The total \textsuperscript{4}He yield was estimated from the measured mass and the helium concentration of 194.2\,nccSTP/mg reported by \citet{farley2000}, a value verified independently in-lab. The atmospheric \textsuperscript{4}He present at the start of the experiment, constrained by the initial permeation steps, was 1.47 $\pm$ 0.18\,nccSTP. The full step-heat dataset is given in the Supplementary Information (Table~\ref{tab:si-durango-stepheat}).

\subsection{Hydrothermal goethite experiment}\label{sec:methods-goethite}

The sample is dense colloform goethite Win-06-01b, collected as float in the Flinders Ranges, South Australia \citep{waltenberg2012,monteiro2017}. It yields one of the oldest known goethite U-Th/He ages, 180--135\,Ma \citep{waltenberg2012}. Win-06-01b was irradiated with 220\,MeV protons for 24 hours at the Francis H. Burr Proton Therapy Center, Massachusetts General Hospital, following \citet{shuster2005} to produce a uniform \textsuperscript{3}He distribution. Aliquots of this same irradiated material were previously step heated in vacuum by \citet{monteiro2017} and \citet{farley2024}, whose results agree well with one another. 

An aliquot of 1.18\,mg Win-06-01 goethite blades shorter than 2\,mm was transferred to the ampoule tip in the presence of an ionizing fan. Less than one microlitre of distilled water was added through a glass capillary, submerging the goethite. The ampoule was tapped to release bubbles before it was flash-frozen in liquid nitrogen, the it was transferred to a $-10$\textdegree C saltwater-ice bath so that the water ice remained solid while atmospheric gases were pumped away. After pumping to <0.1 torr, the saltwater bath was swapped with liquid nitrogen to prevent heating of the sample or melting of the water ice buffer, and the ampoule was fused closed. 

The sealed ampoule, designated HT-Win-06-01b, retained all of its original water through the entire experiment (Fig.~\ref{fig:ampoule-post}), across a pressure contrast of more than eleven orders of magnitude between an internal vapor pressure of roughly 40\,bar at 250\textdegree C \citep{haar1984} and the high vacuum ( 10\textsuperscript{-7} torr) pressures of the lamp cell.

\begin{figure}[htbp]
\centering
\includegraphics[width=0.7\textwidth]{zoomout2.jpg}
\caption{HT-Win-06-01b after completion of the step heats, before breaking for the final furnace step. At room temperature the ampoule is about half gas, one third liquid water, and one sixth goethite shards under 2\,mm long. All of the original water remains, across a pressure contrast sustained for over a week at 250\textdegree C. The gel-like lining on the inner wall is likely amorphous silica dissolved from the ampoule and redeposited by the thermally cycled hydrothermal fluid.}
\label{fig:ampoule-post}
\end{figure}

Before heating, HT-Win-06-01b was held at 70\textdegree C for four days under vacuum to remove atmospheric helium entrained in the ampoule. The experiment then ran for about 40 days over 170--250\textdegree C, with step durations from 2 hours to about a week, through multiple prograde and retrograde cycles designed to test for multiple diffusion domain behavior. We did not heat above 250\textdegree C because of the risk of rupture due to high vapor pressure.

Temperature records show oscillations of about 1\textdegree C amplitude around each set point. While we did not conduct a detailed study to charectarize temperature gradients throughout the ampoule during step heats, we suggest that the apatite experiment of \S\ref{sec:methods-apatite} provides an independent check that recorded temperatures are accurate. 

For the final step, the ampoule was broken and dried overnight at 50\textdegree C, then helium was completely degassed at $\sim$1100\textdegree C in a resistance furnace and measured on a Thermo Fisher Helix SFT mass spectrometer. 

\subsection{Heating and helium isotope analysis}\label{sec:methods-analysis}

Sealed ampoules were washed in ethanol and air dried, then placed in copper foil packets and mounted to a K-type thermocouple in an ultra-high-vacuum projector-lamp heating cell \citep{farley1999}. The lamp cell was pumped to high vacuum and held at 100\textdegree C for 30 minutes to desorb adsorbed species. For encapsulated Durango apatite, less than 0.5\,\textperthousand{} of the radiogenic helium would be expected to escape during this pre-treatment, about 0.01\,nccSTP, using the diffusion parameters of \citet{farley2000} and the forward equations of \citet{fechtigkalbitzer1966}.

Helium isotope abundances for the Win-06-01b goethite were measured on a Thermo Fisher Helix-MC-Plus\textsuperscript{10K}, and 4He abundances were measured with isotope dilution on a Pfeiffer QMS 200 mass spectrometer, both at Caltech. Evolved gases were purified over a liquid-nitrogen charcoal trap, a titanium sponge getter, and a SAES NP10 getter. For the Win-06-01b step heats, \textsuperscript{3}He was counted on a pulse-counting electron multiplier and \textsuperscript{4}He on a Faraday detector. Sensitivities were tracked with measurements of calibrated gas standards every 8 hours, giving internal reproducibilities of 5\% for \textsuperscript{3}He and 2\% for \textsuperscript{4}He (2$\sigma$).

For HT-Win-06-01b, a cold blank gave a negligible \textsuperscript{3}He blank and a \textsuperscript{4}He blank of about 1.25\,pccSTP, comparable to the 4He signals measured during low-temperature permeation steps, but 3He was well-above blank for these steps. Pump-out between steps was 1--5 minutes.

\subsection{Electron microscopy of thermally treated phases}\label{sec:methods-sem}

To assess whether encapsulation preserves reactive phases at experimental temperatures, single crystals of pyrite and arsenopyrite and aggregates of Bahia goethite morphologically similar to Win-06-01b \citep{shustervasconcelos2005,monteiro2017} were subjected to heat treatment in both buffered and unbuffered configurations: unbuffered samples were heated either in vacuum in the lamp-heating apparatus of \citet{farley1999} or in open air while buffered samples were sealed in silica glass ampoules with fused sulfur, mineral realgar, or water for pyrite, arsenopyrite, and goethite respectively. Sulfides were held at 400\textdegree C for at least 4 hours and goethite at 300\textdegree C for 4--11 hours. The ampoules were then broken and the samples air dried, then mounted on carbon tape, carbon coated, and imaged in secondary electron mode on a ZEISS 1550VP field emission SEM at Caltech.

% =============================================================
\section{Permeation Systematics of Fused Silica Ampoules}\label{sec:theory}
% =============================================================

\subsection{Helium budget of a single step}\label{sec:budget}

Consider one heating step. Helium enters the ampoule interior from the sample and leaves it by permeation through the wall, so

\begin{equation}
\label{eq:permlag}
n_{t} = n_{0} + \omega - m .
\end{equation}

Here $n_{0}$ is the helium in the ampoule at the start of the step, $n_{t}$ is the helium remaining at its end, $\omega$ is the helium released by the sample during the step, and $m$ is the helium that permeated out and was measured. Only $m$ is observed. The quantity we want is $\omega$.

Steps are consecutive, so the helium inventory quantified between steps applies to both the previous and the subsequent step:

\begin{equation}
\label{eq:carry}
n_{t}^{\,i} = n_{0}^{\,i+1}, \qquad n_{0}^{\,i} = n_{t}^{\,i-1} .
\end{equation}

Recovering $\omega$ for a step requires knowledge of the ampoule's helium content immediately before and immediately after it. Sections~\ref{sec:permeation-steps} and~\ref{sec:bounding} give the two routes we use to obtain those quantities.

When $\omega$ is large compared with $n_{0}$ and $n_{t}$, Eq.~\eqref{eq:permlag} reduces to $\omega \approx m$ and no correction is needed. That condition does not hold for hour-long steps at 100--400\textdegree C, such as those used on the Durango apatite, but it is approached for long steps used with retentive phases, such as Win-06-01b (\S\ref{sec:bounding}).

\subsection{Assumptions and the permeation rate law}\label{sec:assumptions}

Helium permeation through silica glass is a well-understood first order process \citep{norton1953,altemose1961,swets1961,masaryk1968}. Two conditions must hold before permeation can be described by a first-order rate law.

\begin{enumerate}
\item The helium partial pressure outside the ampoule is negligible compared with that inside. The ampoule and diffusion cell volumes differ by a factor of roughly 10,000 \citep{farley1999}, so this holds throughout.
\item The helium concentration profile across the glass wall is at diffusive steady state. This is established within 30--70 minutes at 80\textdegree C for the ampoules used here, and it is disturbed whenever the temperature changes, because the glass/gas partitioning of helium is temperature dependent. We therefore hold the ampoule at the permeation temperature for at least one diffusive timescale $\tau = L^{2}/D$ after every temperature change before measuring, where $L$ is the wall thickness and $D$ the diffusivity of helium in silica glass \citep{swets1961}.
\end{enumerate}

Given these conditions, permeation is first order in the amount of helium present:

\begin{equation}
\label{eq:permode}
\frac{dn}{dt} = -GKnRT ,
\end{equation}
 
where $n$ is the number of helium atoms, $G$ is a geometric factor with units of cm\textsuperscript{-2}, $K$ is the permeation constant for helium in fused silica, which follows an Arrhenius law $K(T) = K_{0}e^{-E_{a}/RT}$ and has been measured repeatedly \citep{norton1953,altemose1961,swets1961,masaryk1968}, $R$ is the gas constant, and $T$ is temperature in Kelvin. Equation~\eqref{eq:permode} is the standard permeation equation \citep{norton1953,altemose1961,swets1961,masaryk1968}, written in terms of atom counts rather than partial pressures via the ideal gas law and taking the external partial pressure as zero. 

For a cylindrical wall with negligible permeation through the ends, the geometric factor is 

\begin{equation}
\label{eq:geomfactor}
G = \frac{2}{a^{2}\ln(b/a)} ,
\end{equation}

with $a$ the internal and $b$ the external radius \citep{masaryk1968}. There is no dependence on length. $G$ varies with ampoule construction and must be determined for each experiment.

\subsection{Recovering the residual helium with permeation steps}\label{sec:permeation-steps}

At temperatures low enough that the sample releases no helium, $\omega = 0$ and the ampoule's helium inventory changes only by permeation. The measured helium from such a step then is a measurement of the helium inventory of the ampoule. To this end, a useful expression for the fractional yield of an isothermal step is derived, first through integration of Eq.~\eqref{eq:permode} by separation of variables, yielding $n_t = n_0e^{-GKRTt}$, then by relating the measured value to the tube contents by Eq.~\eqref{eq:pemlag} $\left(m(\omega=0)~=~n_0-n_t\right)$.

\begin{equation}
\label{eq:permsol}
^{He}F = \frac{m}{n_{0}} = 1 - e^{-\lambda t} , \qquad \lambda = KGRT ,
\end{equation}

This equation is applied to recover the helium inventory at the beginning of the step.

\begin{equation}
\label{eq:n0perm}
n_{0} = \frac{m}{^{He}F} .
\end{equation}

where $^{He}F$ is the fraction of the ampoule's helium inventory that permeates out during a step of duration $t$, and $\lambda$ is the ampoule's permeation rate constant. Collecting $K$, $G$, $R$, and $T$ into $\lambda$ is useful because $\lambda$ is the quantity that can be measured directly during an experiment (\S\ref{sec:rate-constant}), whereas $K$ and $G$ cannot be separated without independent knowledge of the ampoule geometry.
 
We call such a low-temperature step a \emph{permeation step}. A permeation step does double duty: through Eq.~\eqref{eq:carry} it constrains both the helium remaining after the preceding high-temperature step and the helium inherited by the following one. Placing a permeation step after every high-temperature step therefore supplies $n_{t}$ and $n_{0}$ throughout the experiment, and Eq.~\eqref{eq:permlag} returns $\omega$. This is the approach used to reduce the encapsulated Durango apatite step heat data(\S\ref{sec:results-apatite}).

\subsection{Measuring the ampoule's rate constant during the experiment}\label{sec:rate-constant}

Applying Eq.~\eqref{eq:n0perm} requires $\lambda$. Published permeabilities for anhydrous silica glass \citep{norton1953,altemose1961,swets1961,masaryk1968} combined with a measured $G$ give an estimate, but they should not be relied upon, and they cannot be relied upon at all in hydrothermal experiments: water entering the glass changes its permeability \citep{shelby1972,tsuru2011}, and water-bearing ampoules show anomalously high rate constants (\S\ref{sec:discussion-rate-constant}). The rate constant should therefore be measured on the ampoule actually in use.

Because permeation remains first order regardless of the values of $K$ and $G$, a series of consecutive isothermal permeation steps determines $\lambda$ and $n_{0}$ together with a non-linear least squares solve.

\begin{equation}
\label{eq:incyield}
m_{i+1} = n_{0}\left(1 - e^{-\lambda t_{i+1}}\right) - n_{0}\left(1 - e^{-\lambda t_{i}}\right) .
\end{equation}

We fit individual step yields rather than cumulative yields, because cumulative fitting propagates the errors of early measurements into later measurements, unduly biasing the solver to fit earlier measurements. We used the Levenberg-Marquardt implementation in SciPy \citep{virtanen2020}, weighted by each measurement's own uncertainty, with initial guesses $\lambda \approx 6\times10^{-6}$\,s\textsuperscript{-1} and $n_{0} \approx m_{i}/(1-e^{-\lambda t_{i}})$ evaluated at the first measurement of each series. Uncertainties come from the fit covariance matrices. The fit is poorly conditioned unless enough time has elapsed for the cumulative yield to show curvature, which in practice means letting the permeation rate slow noticeably across three or more consecutive measurements. The exercise can technically be done with as few as two consecutive measurements, although the solver would be unable to estimate uncertainties on the fit parameters. 

\subsection{Bounding the residual helium for long steps}\label{sec:bounding}

Permeation steps require a measurable signal at low temperature. When the sample is retentive and its yields are small, the low-temperature signal can fall to blank level and Eq.~\eqref{eq:n0perm} cannot be applied. This is the situation for \textsuperscript{4}He in the goethite experiment (\S\ref{sec:results-rate-constant}). An alternative is available when steps are long compared with the permeation timescale of the ampoule, which is common for retentive phases.

During a sufficiently long step the ampoule approaches a steady state in which the helium outflux of permeation is balanced by release from the sample, so the residual inventory stops growing while measurable helium continues to accumulate. As time proceeds after this condition is met, the helium inventory becomes relatively small compared to the measured abundance and $\omega \approx m$ becomes an increasinly accurate approximation. Additionaly, an upper bound constraint on the ampoule's helium inventory can be calculated from the high temperature experimental measurement alone. The residual is then bounded above by

\begin{equation}
\label{eq:steadystate}
n^{ss} \leq \frac{m}{\lambda t - 1} ,
\end{equation}

which holds once $\lambda t$ is large enough for steady state to be assumed. When Eq.~\eqref{eq:steadystate} is applied to a long step and to the step preceding it, Eq.~\eqref{eq:permlag}, can be used to bracket $\omega$, the true abundance of helium evolved from the sample:

\begin{equation}
\label{eq:bound}
m_{1} - \frac{m_{0}}{(\lambda t)_{0} - 1} \leq \omega_{1} \leq m_{1}\left(1 + \frac{1}{(\lambda t)_{1} - 1}\right) .
\end{equation}

Equation~\eqref{eq:steadystate} assumes the sample's release rate is constant at a given temperature. In practice the release rate declines with time for volume diffusion \citep{wolf1996,reinersfarley1999,farleymckeon2015,farley2018,farley2024} and for first-order processes \citep{yakubovich2010,wilson1973}, which makes the true residual smaller than the bound and Eq.~\eqref{eq:steadystate} correspondingly conservative. The exceptions are initially inhomogeneous helium distributions at low cumulative yields \citep{goldsmith2020,ketcham2011,shuster2004,farley1996} and radiation damage annealing \citep{cherniak2019,guenthner2013,flowers2007,flowers2009,shuster2006}, which can cause can cause helium release to accelerate.


% =============================================================
\section{Results}\label{sec:results}
% =============================================================

\subsection{Encapsulation stabilizes reactive phases during heating}\label{sec:results-stabilization}

Pyrite, arsenopyrite, and goethite heated without a buffering phase, in vacuum or in open air, develop surface roughening and contraction cracks (Fig.~\ref{fig:stabilization-grid}, right column). The same phases heated at the same temperatures inside sealed silica glass ampoules containing sulfur, realgar, or water retain the smooth surfaces and clear cleavage faces of the untreated controls (Fig.~\ref{fig:stabilization-grid}, centre column). For goethite, cracks in the unbuffered material run both parallel and perpendicular to the long axis of the aggregate. Goethite heated in open air turned red; goethite heated in vacuum remained black.

\begin{figure}[htbp]
\centering
\setlength{\tabcolsep}{4pt}
\begin{tabular}{cccc}
 & \textbf{Control} & \textbf{Buffered} & \textbf{Unbuffered} \\[2pt]

\rotatebox[]{90}{\textbf{Pyrite}} & % origin=c
\includegraphics[width=0.27\textwidth]{LVPY_control_10um0.pdf} &
\includegraphics[width=0.27\textwidth]{LVPy_buff_600_8Hr_10um07.pdf} &
\includegraphics[width=0.27\textwidth]{LVPy_open_400_5Hr_2um11.pdf} \\[4pt]

\rotatebox[]{90}{\textbf{Arsenopyrite}} &
\includegraphics[width=0.27\textwidth]{Asp_control_2um7.pdf} &
\includegraphics[width=0.27\textwidth]{asp_buff_400_8Hr_10um12.pdf} &
\includegraphics[width=0.27\textwidth]{Asp_bigClosed_400C_4hr_10um5.pdf} \\[4pt]

\rotatebox[]{90}{\textbf{Goethite}} & 
\includegraphics[width=0.27\textwidth]{bahiaGoeControl_20um4.pdf} &
\includegraphics[width=0.27\textwidth]{Goeth_Water_300C_4hr_10um6.pdf} &
\includegraphics[width=0.27\textwidth]{Goethite_Vacuum_300C_11hr_10u9.pdf} \\

\end{tabular}
\caption{Secondary electron images of pyrite, arsenopyrite, and goethite before heat treatment (control), after heat treatment buffered within a sealed silica glass ampoule, and after heat treatment unbuffered. Sulfides were heated at 400\textdegree C for at least 4 hours, buffered with fused sulfur (pyrite) or mineral realgar (arsenopyrite); goethite was heated at 300\textdegree C, buffered with water for 4 hours and unbuffered in vacuum for 11 hours. Buffered samples retain the smooth surfaces and clear cleavage faces of the controls, with no evidence of decomposition. Unbuffered samples show contraction cracks and surface roughening. For goethite, cracks in the unbuffered material run both parallel and perpendicular to the long axis of the polycrystalline aggregate.}
\label{fig:stabilization-grid}
\end{figure}

\subsection{Recovering the expected arrhenius plot for encapsulated Durango apatite}\label{sec:results-apatite}

Without correction, apparent diffusivities from the encapsulated Durango apatite step heat are irregular and would not be recognized as belonging to this well-characterized sample (Fig.~\ref{fig:durango-arrhenius}). Prograde sequences give apparent diffusivities that are too low and retrograde sequences give higher values, the signature expected from Eq.~\eqref{eq:permlag}: helium released during a hot step is measured partly during subsequent, cooler steps.

Applying the permeation-step correction of \S\ref{sec:permeation-steps} removes this behavior. The corrected diffusivities agree, within uncertainty, with the Arrhenius relationship reported by \citet{farley2000} for Durango apatite of the same 150\,\textmu m half-prism width, except for the first two plotted points and the 400\textdegree C step (Fig.~\ref{fig:durango-arrhenius}). \citet{farley2000} likewise observed anomalously low apparent diffusivities in early steps. Above 300\textdegree C the corrected array curves downward, a feature documented for Durango apatite and attributed to annealing of radiation damage and crystal defects \citep{wolf1996,farley2000,shuster2006}. Four of the 25 high-temperature steps had permeation lag corrected yields within 2$\sigma$ of zero and are excluded from the figure.

The encapsulated sample reproduces the diffusion behavior measured on unencapsulated material. Permeation lag can be deconvolved from the sample's own helium release.

\begin{figure}[htbp]
\centering
\includegraphics[width=0.75\textwidth]{encapsulated_DApErr.pdf}
\caption{Arrhenius plot of a helium step heat of silica glass encapsulated Durango apatite, with 2$\sigma$ uncertainties. Without permeation lag corrections the apparent diffusivities are highly irregular. With corrections the sample agrees with the diffusion parameters reported by \citet{farley2000} for the appropriate sieve size ($r$ = 150\,\textmu m). Beyond 300\textdegree C the activation energy decreases to give a slight downward curl, likely caused by annealing of crystal defects and radiation damage \citep{wolf1996,farley2000}. Apparent diffusivities are calculated with the equations of \citet{fechtigkalbitzer1966}. Four corrected diffusivities are within 2$\sigma$ of zero and are excluded. The full dataset is given in Table~\ref{tab:si-durango-stepheat}.}
\label{fig:durango-arrhenius}
\end{figure}

\subsection{Rate constant of the hydrothermal goethite ampoule}\label{sec:results-rate-constant}

\textsuperscript{4}He signals during the low-temperature permeation steps of HT-Win-06-01b were at blank level, so \textsuperscript{4}He could not be treated by the permeation-step route of \S\ref{sec:permeation-steps}; however, \textsuperscript{3}He signals were well above blank, permitting in-situ measurement of the rate constant of helium permeation. Ten series of three to four consecutive isothermal permeation steps, held at 90--110\textdegree C for 120--180 minutes each with 26 minutes of cooling before inlet, were fit to Eq.~\eqref{eq:incyield} (Fig.~\ref{fig:permeation-fits}). Each series is well described by the first-order model, which supports the rate law of Eq.~\eqref{eq:permode} for this ampoule and yields $\lambda$ directly. One further series that otherwise qualified was poorly fit with high uncertainties and was excluded from this analysis.

\begin{figure}[htbp]
\centering
\includegraphics[width=0.8\textwidth]{NaiveIsothermalFits_Grid.pdf}
\caption{Cumulative \textsuperscript{3}He yield against cumulative time for each series of consecutive isothermal permeation steps on HT-Win-06-01b, with the non-linear least squares fit of $\lambda$ and $n_{0}$ to Eq.~\eqref{eq:incyield} overlain. Each series is held at a single temperature between 90 and 110\textdegree C. The curvature of each series is what conditions the fit: a series in which the permeation rate has not yet slowed appreciably constrains $\lambda$ and $n_{0}$ only weakly (\S\ref{sec:rate-constant}). The rate constants recovered here are converted to permeabilities, $K = \lambda/GRT$, and plotted against temperature in Fig.~\ref{fig:ampoule-permeability}.}
\label{fig:permeation-fits}
\end{figure}

Converting these rate constants to permeabilities through $K = \lambda/GRT$, every apparent \textsuperscript{3}He permeability equals or exceeds the anhydrous literature value (Fig.~\ref{fig:ampoule-permeability}). A regression of $\ln K$ against inverse temperature over 90--110\textdegree C gives $R^{2} = \ChTwoPermRtwo$ across \ChTwoPermN\ series, with an activation energy of \ChTwoPermEa\ $\pm$ \ChTwoPermEaErr\,cal/mol (2 s.e.), significantly higher than the value expected for silica glass without water \citep{swets1961,jones1953}, closer to glasses rich in network-modifying cations \citep{norton1953,altemose1961}. 

We tested whether the elevated apparent permeability could be an artifact of continued permeation during the cooling interval before each measurement. A finite-difference model of the ampoule, using the measured cooling schedule, predicts that excess helium permeated during the cooling period would produce an offset that is nearly constant, weakly dependent on temperature, and much smaller than the difference between the measured and literature permeability values (Fig.~\ref{fig:ampoule-permeability}). The elevated rate constant is not a cooling artifact. 

\begin{figure}[htbp]
\centering
\includegraphics[width=0.8\textwidth]{NaiveLambda_Arrhenius.pdf}
\caption{Apparent permeation coefficients for HT-Win-06-01b against inverse temperature, converted from the rate constants fitted in Fig.~\ref{fig:permeation-fits} via $K = \lambda/GRT$. The activation energy is higher than for anhydrous silica glass, as expected \citep{masaryk1968,shelby1972}. The apparent \textsuperscript{3}He permeability exceeding the literature values of \citet{swets1961} and \citet{jones1953} is not expected, and is attributed to displacement of gas-phase volume by liquid water (\S\ref{sec:discussion-rate-constant}). Points labelled ``Diffusion Model with cooling, Anhydrous'' give the excess apparent permeability predicted for a dry ampoule that continued to permeate during the 26 minute cooldown before measurement, simulated over the same 120--180 minute holds at 90--110\textdegree C. Their scatter reflects differences in hold duration and is small; the offset from the literature values is near constant and much smaller than the uncertainty on the fitted values.}
\label{fig:ampoule-permeability}
\end{figure}

Because the measured permeabilities and their activation energy exceed the anhydrous literature permeabilities \citep{swets1961,altemose1961,jones1953}, the literature-implied rate constants can safely be taken as lower bound on $\lambda$ for this ampoule, and Eq.~\eqref{eq:steadystate} evaluated with it returns a conservative upper bound on the residual helium. We therefore reduced HT-Win-06-01b by the bounding approach of \S\ref{sec:bounding}, taking $\lambda$ from the \textsuperscript{4}He permeability of \citet{swets1961} and, for \textsuperscript{3}He, from that value scaled by the ratio of \citet{jones1953}.

The first three experimental steps have $\lambda t < 4$ and no bound could be evaluated; they are excluded from further analysis. Later steps with $\lambda t < 4$ could be treated, because bounds on the preceding steps' residuals were available. Tables~\ref{tab:winsor-3he} and~\ref{tab:winsor-4he} give the resulting dataset.

\clearpage
\begin{landscape}
\begin{footnotesize}
\begin{longtable}{r r r r r r r r r r}
\caption{$^3$He step-heat dataset for HT-Win-06-01b under the ``treatment as uncertainty'' method: measured and cumulative $^3$He, the upper-bound residual ampoule content ($n_{t,\mathrm{UB}}$), the lower- and upper-bound true step production ($\omega_{\mathrm{LB}}$, $\omega_{\mathrm{UB}}$), and the corresponding lower-bound, central, and upper-bound apparent diffusivities. ``--'' denotes early steps with $\lambda t < 4$, for which bounds could not be determined (see \S\ref{sec:rate-constant}). Six rows have derived values marked with ``*'': these steps have measured $^3$He $<0.1$ pccSTP and are excluded as low-signal from all plots (Fig.~\ref{fig:ratio-spectra} and others derived from this dataset).}\label{tab:winsor-3he}\\
\toprule
Temp & Time & $^3$He meas. & Cum. $^3$He & $n_{t,\mathrm{UB}}$ & $\omega_{\mathrm{LB}}$ & $\omega_{\mathrm{UB}}$ & $\ln(D/a^2)_{\mathrm{LB}}$ & $\ln(D/a^2)_{m}$ & $\ln(D/a^2)_{\mathrm{UB}}$ \\
(°C) & (min) & (pccSTP) & (pccSTP) & (pccSTP) & (pccSTP) & (pccSTP) &  &  &  \\
\midrule
\endfirsthead
\multicolumn{10}{l}{\textit{(continued)}} \\
\toprule
Temp & Time & $^3$He meas. & Cum. $^3$He & $n_{t,\mathrm{UB}}$ & $\omega_{\mathrm{LB}}$ & $\omega_{\mathrm{UB}}$ & $\ln(D/a^2)_{\mathrm{LB}}$ & $\ln(D/a^2)_{m}$ & $\ln(D/a^2)_{\mathrm{UB}}$ \\
(°C) & (min) & (pccSTP) & (pccSTP) & (pccSTP) & (pccSTP) & (pccSTP) &  &  &  \\
\midrule
\endhead
\midrule
\multicolumn{10}{r}{\textit{continued on next page}} \\
\endfoot
\bottomrule
\endlastfoot
170 & 300 & 0.9630 & 0.0000 & 1.1660 & -- & -- & -- & -- & -- \\
180 & 463 & 0.5442 & 0.9630 & 0.2413 & -- & -- & -- & -- & -- \\
190 & 550 & 0.5292 & 1.6241 & 0.1538 & -- & -- & -- & -- & -- \\
200 & 500 & 0.5190 & 2.2112 & 0.1430 & 0.3651 & 0.6619 & -22.158 & -21.775 & -21.503 \\
210 & 475 & 0.5770 & 2.7646 & 0.1440 & 0.4341 & 0.7210 & -21.715 & -21.406 & -21.160 \\
220 & 450 & 0.6333 & 3.3799 & 0.1447 & 0.4894 & 0.7781 & -21.345 & -21.068 & -20.843 \\
230 & 425 & 0.7899 & 4.0551 & 0.1685 & 0.6452 & 0.9584 & -20.823 & -20.604 & -20.392 \\
225 & 908 & 0.7513 & 4.8885 & 0.0716 & 0.5827 & 0.8229 & -21.516 & -21.246 & -21.148 \\
215 & 450 & 0.1415 & 5.6842 & 0.0349 & 0.0699 & 0.1764 & -22.836 & -22.124 & -21.900 \\
205 & 520 & 0.0763 & 5.8529 & 0.0184$^*$ & 0.0413$^*$ & 0.0947$^*$ & -23.479$^*$ & -22.863$^*$ & -22.646$^*$ \\
195 & 1106 & 0.0670 & 5.9491 & 0.0080$^*$ & 0.0487$^*$ & 0.0750$^*$ & -24.053$^*$ & -23.731$^*$ & -23.618$^*$ \\
185 & 740 & 0.0227 & 6.0246 & 0.0049$^*$ & 0.0147$^*$ & 0.0276$^*$ & -24.839$^*$ & -24.404$^*$ & -24.207$^*$ \\
175 & 920 & 0.0180 & 6.0610 & 0.0036$^*$ & 0.0131$^*$ & 0.0216$^*$ & -25.168$^*$ & -24.847$^*$ & -24.665$^*$ \\
200 & 260 & 0.0224 & 6.0814 & 0.0163$^*$ & 0.0189$^*$ & 0.0388$^*$ & -23.534$^*$ & -23.360$^*$ & -22.812$^*$ \\
210 & 520 & 0.0812 & 6.1081 & 0.0181$^*$ & 0.0649$^*$ & 0.0993$^*$ & -22.984$^*$ & -22.758$^*$ & -22.554$^*$ \\
220 & 1167 & 0.5766 & 6.1914 & 0.0447 & 0.5585 & 0.6213 & -21.587 & -21.553 & -21.475 \\
230 & 370 & 0.3577 & 6.7796 & 0.0903 & 0.3130 & 0.4480 & -20.948 & -20.811 & -20.579 \\
240 & 320 & 0.6767 & 7.1563 & 0.1776 & 0.5864 & 0.8543 & -20.062 & -19.913 & -19.667 \\
245 & 250 & 0.5549 & 7.8582 & 0.1888 & 0.3773 & 0.7437 & -20.176 & -19.778 & -19.474 \\
250 & 210 & 0.6934 & 8.4743 & 0.2786 & 0.5047 & 0.9721 & -19.626 & -19.296 & -18.942 \\
245 & 270 & 0.4659 & 9.2086 & 0.1413 & 0.1872 & 0.6071 & -20.801 & -19.874 & -19.601 \\
240 & 320 & 0.2857 & 9.7206 & 0.0749 & 0.1445 & 0.3606 & -21.176 & -20.487 & -20.250 \\
235 & 370 & 0.2199 & 10.0357 & 0.0517 & 0.1450 & 0.2716 & -21.284 & -20.864 & -20.650 \\
230 & 420 & 0.1400 & 10.2736 & 0.0303 & 0.0883 & 0.1703 & -21.885 & -21.421 & -21.224 \\
225 & 652 & 0.1494 & 10.4279 & 0.0206 & 0.1191 & 0.1700 & -22.009 & -21.781 & -21.650 \\
230 & 500 & 0.1550 & 10.5908 & 0.0271 & 0.1344 & 0.1820 & -21.605 & -21.462 & -21.299 \\
235 & 450 & 0.1951 & 10.7582 & 0.0361 & 0.1681 & 0.2313 & -21.258 & -21.107 & -20.936 \\
240 & 400 & 0.2425 & 10.9638 & 0.0483 & 0.2063 & 0.2908 & -20.913 & -20.750 & -20.566 \\
245 & 350 & 0.3150 & 11.2328 & 0.0690 & 0.2667 & 0.3840 & -20.495 & -20.326 & -20.125 \\
250 & 300 & 0.3716 & 11.5689 & 0.0939 & 0.3026 & 0.4656 & -20.181 & -19.973 & -19.743 \\
250 & 300 & 0.3640 & 11.9627 & 0.0899 & 0.2700 & 0.4539 & -20.261 & -19.959 & -19.734 \\
250 & 300 & 0.3452 & 12.3490 & 0.0852 & 0.2553 & 0.4304 & -20.285 & -19.979 & -19.755 \\
250 & 10250 & 8.4829 & 12.7167 & 0.0494 & 8.3977 & 8.5324 & -19.989 & -19.976 & -19.969 \\
240 & 785 & 0.2261 & 21.2187 & 0.0209 & 0.1767 & 0.2470 & -21.026 & -20.778 & -20.689 \\
\midrule
\textbf{Total} &  &  & 146.1483 &  &  &  &  &  &  \\
\end{longtable}
\end{footnotesize}
\end{landscape}

\begin{landscape}
\begin{footnotesize}
\begin{longtable}{r r r r r r r r r r r r r}
\caption{$^4$He step-heat dataset and absolute $^4$He/$^3$He ratios for HT-Win-06-01b under the ``treatment as uncertainty'' method, structured as in Table~\ref{tab:winsor-3he}. ``--'' denotes steps for which bounds could not be determined. The Total row reports the bulk (total-yield) $^4$He/$^3$He ratio, i.e. the normalizing value used to bulk-normalize these ratios for the plot in Fig.~\ref{fig:ratio-spectra}. Six rows have derived values marked with ``*'': these steps have measured $^3$He $<0.1$ pccSTP and are excluded as low-signal from all plots, as in Table~\ref{tab:winsor-3he}.}\label{tab:winsor-4he}\\
\toprule
Temp & Time & $^4$He meas. & Cum. $^4$He & $n_{t,\mathrm{UB}}$ & $\omega_{\mathrm{LB}}$ & $\omega_{\mathrm{UB}}$ & $\ln(D/a^2)_{\mathrm{LB}}$ & $\ln(D/a^2)_{m}$ & $\ln(D/a^2)_{\mathrm{UB}}$ & $(^4\mathrm{He}/^3\mathrm{He})_{\mathrm{LB}}$ & $(^4\mathrm{He}/^3\mathrm{He})_{m}$ & $(^4\mathrm{He}/^3\mathrm{He})_{\mathrm{UB}}$ \\
(°C) & (min) & (nccSTP) & (nccSTP) & (nccSTP) & (nccSTP) & (nccSTP) &  &  &  &  &  &  \\
\midrule
\endfirsthead
\multicolumn{13}{l}{\textit{(continued)}} \\
\toprule
Temp & Time & $^4$He meas. & Cum. $^4$He & $n_{t,\mathrm{UB}}$ & $\omega_{\mathrm{LB}}$ & $\omega_{\mathrm{UB}}$ & $\ln(D/a^2)_{\mathrm{LB}}$ & $\ln(D/a^2)_{m}$ & $\ln(D/a^2)_{\mathrm{UB}}$ & $(^4\mathrm{He}/^3\mathrm{He})_{\mathrm{LB}}$ & $(^4\mathrm{He}/^3\mathrm{He})_{m}$ & $(^4\mathrm{He}/^3\mathrm{He})_{\mathrm{UB}}$ \\
(°C) & (min) & (nccSTP) & (nccSTP) & (nccSTP) & (nccSTP) & (nccSTP) &  &  &  &  &  &  \\
\midrule
\endhead
\midrule
\multicolumn{13}{r}{\textit{continued on next page}} \\
\endfoot
\bottomrule
\endlastfoot
170 & 300 & 0.0192 & 0.0000 & 0.0280 & -- & -- & -- & -- & -- & -- & -- & -- \\
180 & 463 & 0.0305 & 0.0192 & 0.0152 & -- & -- & -- & -- & -- & -- & -- & -- \\
190 & 550 & 0.0558 & 0.0660 & 0.0179 & -- & -- & -- & -- & -- & -- & -- & -- \\
200 & 500 & 0.1024 & 0.1375 & 0.0313 & 0.0845 & 0.1337 & -26.207 & -25.965 & -25.619 & -- & 197.3 & -- \\
210 & 475 & 0.1460 & 0.2500 & 0.0404 & 0.1147 & 0.1864 & -25.314 & -25.022 & -24.717 & 248.2 & 253.0 & 268.8 \\
220 & 450 & 0.2091 & 0.4055 & 0.0530 & 0.1686 & 0.2620 & -24.407 & -24.152 & -23.876 & 326.3 & 330.1 & 349.3 \\
230 & 425 & 0.3243 & 0.6319 & 0.0760 & 0.2714 & 0.4004 & -23.425 & -23.213 & -22.956 & 403.7 & 410.6 & 432.3 \\
225 & 908 & 0.3483 & 0.9799 & 0.0363 & 0.2722 & 0.3845 & -23.807 & -23.527 & -23.413 & 442.6 & 463.6 & 474.4 \\
215 & 450 & 0.0716 & 1.3487 & 0.0196 & 0.0353 & 0.0912 & -24.945 & -24.225 & -23.976 & 429.1 & 505.9 & 537.1 \\
205 & 520 & 0.0385 & 1.4359 & 0.0103$^*$ & 0.0189$^*$ & 0.0487$^*$ & -25.661$^*$ & -24.942$^*$ & -24.701$^*$ & 392.4$^*$ & 504.2$^*$ & 534.6$^*$ \\
195 & 1106 & 0.0336 & 1.4912 & 0.0043$^*$ & 0.0233$^*$ & 0.0379$^*$ & -26.165$^*$ & -25.796$^*$ & -25.675$^*$ & 446.7$^*$ & 500.9$^*$ & 515.3$^*$ \\
185 & 740 & 0.0079 & 1.5467 & 0.0019$^*$ & 0.0037$^*$ & 0.0098$^*$ & -27.586$^*$ & -26.808$^*$ & -26.592$^*$ & 224.7$^*$ & 350.1$^*$ & 369.1$^*$ \\
175 & 920 & 0.0152 & 1.6020 & 0.0033$^*$ & 0.0133$^*$ & 0.0185$^*$ & -26.475$^*$ & -26.340$^*$ & -26.141$^*$ & 843.2$^*$ & 843.2$^*$ & 982.0$^*$ \\
200 & 260 & 0.0107 & 1.6351 & 0.0091$^*$ & 0.0073$^*$ & 0.0198$^*$ & -25.787$^*$ & -25.411$^*$ & -24.791$^*$ & 375.6$^*$ & 475.5$^*$ & 554.7$^*$ \\
210 & 520 & 0.2166 & 1.6547 & 0.0535$^*$ & 0.2075$^*$ & 0.2701$^*$ & -23.067$^*$ & -23.021$^*$ & -22.786$^*$ & 2667.1$^*$ & 2667.1$^*$ & 3103.8$^*$ \\
220 & 1167 & 0.3187 & 1.8714 & 0.0269 & 0.2652 & 0.3456 & -23.499 & -23.302 & -23.214 & 471.7 & 552.7 & 563.2 \\
230 & 370 & 0.2113 & 2.6766 & 0.0590 & 0.1844 & 0.2703 & -22.391 & -22.249 & -21.993 & 575.1 & 590.8 & 628.0 \\
240 & 320 & 0.4286 & 3.1616 & 0.1243 & 0.3697 & 0.5529 & -21.360 & -21.203 & -20.931 & 612.1 & 633.4 & 674.8 \\
245 & 250 & 0.3988 & 3.6430 & 0.1521 & 0.2746 & 0.5510 & -21.289 & -20.899 & -20.556 & 665.9 & 718.7 & 778.4 \\
250 & 210 & 0.4440 & 4.2359 & 0.2020 & 0.2918 & 0.6460 & -20.906 & -20.469 & -20.072 & 541.7 & 640.2 & 702.8 \\
245 & 270 & 0.3123 & 4.7263 & 0.1053 & 0.1103 & 0.4176 & -22.043 & -20.981 & -20.680 & 469.1 & 670.4 & 721.4 \\
240 & 320 & 0.1927 & 5.0846 & 0.0559 & 0.0874 & 0.2485 & -22.376 & -21.575 & -21.315 & 505.5 & 674.3 & 718.4 \\
235 & 370 & 0.1517 & 5.3031 & 0.0393 & 0.0958 & 0.1910 & -22.386 & -21.922 & -21.688 & 599.8 & 689.6 & 729.9 \\
230 & 420 & 0.0991 & 5.4714 & 0.0236 & 0.0598 & 0.1227 & -22.956 & -22.448 & -22.232 & 608.3 & 707.9 & 745.9 \\
225 & 652 & 0.0995 & 5.5835 & 0.0151 & 0.0760 & 0.1146 & -23.136 & -22.863 & -22.721 & 608.2 & 666.5 & 689.2 \\
230 & 500 & 0.1051 & 5.6969 & 0.0202 & 0.0900 & 0.1253 & -22.679 & -22.523 & -22.345 & 652.1 & 678.1 & 707.5 \\
235 & 450 & 0.1369 & 5.8251 & 0.0279 & 0.1167 & 0.1648 & -22.290 & -22.128 & -21.940 & 674.6 & 701.7 & 734.0 \\
240 & 400 & 0.1781 & 5.9762 & 0.0390 & 0.1502 & 0.2171 & -21.892 & -21.719 & -21.518 & 705.1 & 734.6 & 770.9 \\
245 & 350 & 0.2305 & 6.1807 & 0.0556 & 0.1915 & 0.2861 & -21.479 & -21.290 & -21.070 & 694.4 & 731.9 & 771.7 \\
250 & 300 & 0.2777 & 6.4385 & 0.0758 & 0.2221 & 0.3535 & -21.134 & -20.906 & -20.659 & 703.3 & 747.3 & 793.3 \\
250 & 300 & 0.2654 & 6.7374 & 0.0725 & 0.1896 & 0.3379 & -21.250 & -20.908 & -20.661 & 653.7 & 729.2 & 774.1 \\
250 & 300 & 0.2489 & 7.0231 & 0.0680 & 0.1764 & 0.3169 & -21.282 & -20.932 & -20.647 & 646.8 & 721.0 & 765.4 \\
250 & 10250 & 6.8042 & 7.2901 & 0.0430 & 6.7362 & 6.8472 & -20.704 & -20.691 & -20.682 & 800.6 & 802.1 & 803.3 \\
240 & 785 & 0.1853 & 14.1156 & 0.0187 & 0.1423 & 0.2040 & -21.688 & -21.422 & -21.325 & 770.5 & 819.4 & 838.0 \\
\midrule
\textbf{Total} &  &  & 145.8145 &  &  &  &  &  &  &  & 997.7 &  \\
\end{longtable}
\end{footnotesize}
\end{landscape}

\clearpage

\subsection{Helium release from goethite under hydrothermal conditions}\label{sec:results-goethite}

The ampoule retained its full water charge for the duration of the experiment (Fig.~\ref{fig:ampoule-post}), so hydrothermal conditions were maintained throughout.

Apparent diffusivities for HT-Win-06-01b are lower than those measured for the same sample in vacuum by \citet{monteiro2017} and \citet{farley2024} at every step. This gap widens to approximately seven natural log units after the first few steps (Figs.~\ref{fig:arrhenius-full},~\ref{fig:arrhenius-zoom}).

The \textsuperscript{3}He Arrhenius array sharply declines along a concave-up trend to about 3.5\% cumulative \textsuperscript{3}He yield, then decreases slowly along sub-parallel arrays through the prograde and retrograde cycles between 3.5 and 7\% cumulative yield. Beyond about 7\% cumulative yield the array is linear. A regression over the remaining \ChTwoN\ steps, spanning 7--15\% cumulative \textsuperscript{3}He yield, gives $R^{2} = \ChTwoRtwo$ and an activation energy of \ChTwoEa\ $\pm$ \ChTwoEaErr\,kJ/mol. Treating this array as a single diffusion domain and considering only the uncertainty in the intercept gives a closure temperature of \ChTwoTc\ $\pm$ \ChTwoTcErr\,\textdegree C at a cooling rate of 10\textdegree C/Ma. Total \textsuperscript{3}He released over the experiment was 14.7\% of the sample's total \textsuperscript{3}helium inventory.

The \textsuperscript{4}He array differs in shape. It is concave-down to about 0.1\% cumulative \textsuperscript{4}He yield, after which it follows a roughly linear array about one log unit below the \textsuperscript{3}He array.

On a delta plot constructed against the \textsuperscript{3}He regression as reference line, \textsuperscript{3}He delta values decrease monotonically until about 7\% cumulative loss and are within error of the reference line thereafter (Fig.~\ref{fig:delta-zoom}). \textsuperscript{4}He delta values are more uncertain; they increase initially and settle into a slow increase or plateau about one log unit below the \textsuperscript{3}He reference line by about 5\% cumulative \textsuperscript{3}He yield. A group of low-yield measurements on the first retrograde cycle produce scattered, high-uncertainty delta values just after 4\% cumulative loss.

Bulk-normalized \textsuperscript{4}He/\textsuperscript{3}He ratios begin at very low values and increase steeply, with the rate of increase slowing as the experiment proceeds (Fig.~\ref{fig:ratio-spectra}). At any given cumulative \textsuperscript{3}He yield the hydrothermal experiment records lower bulk-normalized ratios than either vacuum experiment on the same sample, and its approach toward a plateau is the slowest of the three. The hydrothermal experiment did not reach a plateau. The elevated first-step ratio seen in both vacuum experiments is absent here.

\begin{figure}[htbp]
\centering
\includegraphics[width=0.8\textwidth]{winsor_goethite_Diffusivities_Arrhenius.pdf}
\caption{Arrhenius plot of \textsuperscript{3}He and \textsuperscript{4}He apparent diffusivities from HT-Win-06-01b, with step heats of the same Win-06-01b material in vacuum from \citet{monteiro2017} and \citet{farley2024}. The Stage 2 reference line is fit to the step heats of \citet{monteiro2017}. Error bars on the data from this study are omitted for clarity. A diffusion domain corresponding to the Stage 2 reference line would have a closure temperature of about 45\textdegree C at a cooling rate of 10\textdegree C/Ma.}
\label{fig:arrhenius-full}
\end{figure}

\begin{figure}[htbp]
\centering
\includegraphics[width=0.8\textwidth]{winsor_goethite_Diffusivities_Arrhenius_Zoomin.pdf}
\caption{Linear regression of the \textsuperscript{3}He apparent diffusivities of HT-Win-06-01b above 7\% cumulative yield, defining the Stage 1 reference line. A diffusion domain corresponding to this line has a closure temperature of \ChTwoTc\ $\pm$ \ChTwoTcErr\,\textdegree C at a cooling rate of 10\textdegree C/Ma. Apparent diffusivities calculated from \textsuperscript{4}He are sometimes within error of those calculated from \textsuperscript{3}He, but are occasionally resolved at the 2$\sigma$ level and lie about one log unit lower after a cumulative \textsuperscript{3}He yield of 14\%.}
\label{fig:arrhenius-zoom}
\end{figure}

\begin{figure}[htbp]
\centering
\includegraphics[width=0.8\textwidth]{DeltasWinsorGTzoomin_VacuumHydrothermal.pdf}
\caption{Delta plot after \citet{farley2018}, comparable to the $r/r_{0}$ plots of \citet{lovera1991} and \citet{farleyflowers2012}. Delta is the difference between apparent log diffusivity and the Stage 1 reference line, which removes the temperature dependence of the dominant migration process to show how apparent diffusivities evolve through the experiment. \textsuperscript{3}He decreases monotonically, settles on a temporary plateau from 4--6\% fractional yield, then follows the reference line from 7--15\%. \textsuperscript{4}He delta values are more uncertain; they increase initially and settle about one log unit below the \textsuperscript{3}He reference line by about 5\% \textsuperscript{3}He yield.}
\label{fig:delta-zoom}
\end{figure}

\begin{figure}[htbp]
\centering
\includegraphics[width=0.8\textwidth]{Comparative_Ratios_Dtplot_WinsorGT_VacuumHDMKAF.pdf}
\caption{Bulk-normalized \textsuperscript{4}He/\textsuperscript{3}He ratios for Win-06-01b step heats against cumulative fractional \textsuperscript{3}He yield. The hydrothermal experiment records lower ratios at a given yield than either vacuum experiment for the entire duration, apart from the first step. All three begin at very small values and increase steeply, with the rate of increase slowing between 0.05 and 0.2 cumulative yield. The hydrothermal experiment is plotted as an envelope spanning the lower and upper bounds of Eq.~\eqref{eq:bound}. Ten of 34 analyses are excluded, either because $\lambda t$ was too small to evaluate bounds or because the measured \textsuperscript{3}He signal was below 0.1\,pccSTP.}
\label{fig:ratio-spectra}
\end{figure}

% =============================================================
\section{Discussion}\label{sec:discussion}
% =============================================================

\subsection{Elevated rate constants in water-bearing ampoules}\label{sec:discussion-rate-constant}

A rate constant above the anhydrous value is the opposite of the expected result. Water and hydroxyl dissolved in silica glass are reported to decrease helium permeability at all temperatures \citep{tsuru2011,shelby1972,masaryk1968}, on the interpretation that species such as OH\textsuperscript{-} occupy interstitial space in the glass network \citep{norton1953,shackelford1999}. We have observed anomalously high rate constants in water-bearing ampoules consistently, which is what makes in-situ measurement necessary rather than merely advisable.

The likely explanation does not involve the glass. Helium is only sparingly soluble in water and diffuses through it rapidly, roughly $7 \times 10^{-5}$\,cm\textsuperscript{2}/s at room temperature \citep{verhallen1984} against about $2 \times 10^{-8}$\,cm\textsuperscript{2}/s in silica glass \citep{swets1961}. On the timescale of permeation through the wall, transport through the water is effectively instantaneous, so the gas phase stays in equilibrium with the glass over the whole inner surface, including where liquid water contacts it. Adding water therefore reduces the volume available to the gas phase without reducing the area available for permeation. Helium partitions strongly into the gas phase over both water and glass: concentrations in the glass and in the water are about 3\% and 1\% respectively of the gas-phase concentration, using solubilities from \citet{sander2015} and \citet{swets1961}. Displacing gas-phase volume therefore raises the helium partial pressure, and the apparent permeability scales approximately inversely with the gas-phase volume fraction,

\begin{equation}
\label{eq:keff}
K_{eff} = K_{dry}\,\frac{V_{water} + V_{gas}}{V_{gas}} .
\end{equation}

The effect is geometric. The ampoule's helium is concentrated into a smaller volume than the nominal internal volume assumed in calculating $G$, so the discrepancy appears in $K$ when $\lambda$ is decomposed using the nominal geometry. Any condensed phase does this, but the magnitude scales with the volume fraction displaced: water occupied roughly half the internal volume of HT-Win-06-01b, whereas a milligram-scale mineral charge occupies a few percent, which is why the effect is resolvable here and not for encapsulated Durango apatite.

This anomaly is what licenses the reduction used for the goethite dataset. Because the ampoule's own \textsuperscript{3}He measurements show its rate constant to exceed the anhydrous value (Fig.~\ref{fig:ampoule-permeability}), the anhydrous value is a lower bound on $\lambda$, and Eq.~\eqref{eq:steadystate} evaluated with it returns an upper bound on the residual helium and hence a conservative bound on the correction. Had water-bearing ampoules permeated more slowly than anhydrous glass, as the glass-hydration literature predicts, this reduction would not have been available.

The same reasoning constrains the isotope ratios. Isotopic fractionation between measured and true release grows with the \textsuperscript{3}He/\textsuperscript{4}He permeability contrast and with the strength of the lag, so we take the \textsuperscript{4}He permeability to be that of \citet{swets1961}, which the experiment's own \textsuperscript{3}He measurements show to be an underestimate, and the relative \textsuperscript{3}He/\textsuperscript{4}He permeability to be 1.29. That value exceeds the 1.06--1.11 expected from \citet{jones1953} over this temperature range. Both choices widen the reported envelope rather than narrowing it.

\subsection{Helium release from goethite is consistent with multiple diffusion domain systematics}\label{sec:discussion-mdd}

Natural goethite occurs as aggregates of crystallites typically smaller than 1\,\textmu m, and all published goethite step heats are on polycrystalline material of unclear and probably variable crystallite size \citep{monteiro2017,farley2024}. Such a sample should show multiple diffusion domain behavior, and neglecting it produces inaccurate activation energies and frequency factors \citep{farley2018}.

In a multiple diffusion domain system with an initially uniform diffusant, apparent bulk diffusivity at fixed temperature decreases as the least retentive domains are exhausted, and delta values decrease correspondingly, provided the reference line's activation energy matches that of the controlling process \citep{lovera1991,farley2018}. Using the activation energy retrieved here, the \textsuperscript{3}He delta values decrease monotonically (Fig.~\ref{fig:delta-zoom}), which is the expected behavior.

The plateau in the delta plot and the quality of the regression over 7--15\% cumulative \textsuperscript{3}He yield ($R^{2} = \ChTwoRtwo$) suggest that a single domain dominates the frequency factor over this interval, rather than a mixture of domains as is typically observed for polycrystalline hematite \citep{farley2018}. The \textsuperscript{4}He apparent diffusivities are lower than the \textsuperscript{3}He values and increase early in the experiment, which is what rounding of the radiogenic \textsuperscript{4}He profile by natural loss from the least retentive domains would produce \citep{shuster2004,shustervasconcelos2005}.

Vacuum experiments on goethite also record declining delta values at low \textsuperscript{3}He yields that are not mirrored in the \textsuperscript{4}He data \citep{shustervasconcelos2005,farley2024}. One interpretation is migration out of fast pathways along grain boundaries that would retain little helium in nature \citep{shustervasconcelos2005}. Goethite aggregates contain interstices filled with chemisorbed non-stoichiometric water, which is the first water lost in vacuum \citep{naono1987}, and non-stoichiometric hydrogen persists in goethite above 140\textdegree C \citep{yapppoths1995}. Proton-induced \textsuperscript{3}He sited in such interstices may be released with this early water, which would correspond to the low retention domain of \citet{shustervasconcelos2005} and the Stage 1 release of \citet{farley2024}. The wide variability in the size of this early fraction is consistent with the strong dependence of goethite crystallinity, surface area, and non-stoichiometric water content on formation temperature \citep{schwertmann1985}.

These data do not bear on whether helium diffusion in crystalline goethite is anisotropic or modulated by radiation damage \citep{bassal2022}.

\subsection{Dehydration and recrystallization}\label{sec:discussion-alteration}

Bulk goethite is unstable with respect to hematite and water at the pressures and temperatures of this experiment \citep{majzlan2003}, nanoparticles excepted \citep{guobarnard2011,navrotsky2008}. Dehydration or recrystallization could therefore contribute to the observed release, and neither can be excluded.

The electron microscopy of \S\ref{sec:results-stabilization} bounds pervasive dehydration. Goethite heated in vacuum or open air develops contraction cracks through grain interiors and across grain surfaces, while goethite heated under water-saturated conditions does not (Fig.~\ref{fig:stabilization-grid}). This excludes dehydration at the scale that produces visible cracking. It does not exclude dehydration below that scale, and work on the textural evolution of goethite aggregates at low dehydration fractions would be needed to resolve how much can occur before cracks appear.

Recrystallization is harder to bound. Poorly crystalline synthetic goethite recrystallizes under hydrothermal conditions in as little as 16 hours at 125\textdegree C \citep{schwertmann1985}, and our conditions were considerably harsher. The silica dissolved from the ampoule wall and redeposited as a gel-like lining (Fig.~\ref{fig:ampoule-post}) shows that the fluid was capable of dissolving and reprecipitating rock-forming oxides. Whether it recrystallized Win-06-01b specifically is unclear; this sample is among the most crystalline and most helium-retentive goethites yet dated \citep{waltenberg2012,farley2024}, whereas the less crystalline goethites commonly recrystallized in nature appear to lose their helium and reset \citep{monteiro2018,monteiro2014,heim2006}.

The interpretive consequence is the same either way. If dehydration or recrystallization contributed to the release, the apparent diffusivities reported here are upper bounds on the true helium diffusivity in goethite at these temperatures and yields, and the retentivity we infer is a lower bound. Both would then imply that dehydration and recrystallization, rather than volume diffusion, set the pace of helium loss from goethite in nature.

\subsection{Helium retentivity of goethite}\label{sec:discussion-retentivity}

Apparent diffusivities for HT-Win-06-01b are resolved from, and about seven natural log units below, those measured for the same sample in vacuum \citep{monteiro2017,farley2024}. We consider what this implies through closure temperature, isothermal retention, and comparison with other polycrystalline samples. Each applies strictly only to loss by thermally activated volume diffusion from domains of known size, so each carries the caveats of \S\ref{sec:discussion-alteration}, and none permits estimating the diffusivity at infinite temperature.

Treating the Stage 1 regression as a real diffusion domain gives a closure temperature of \ChTwoTc\ $\pm$ \ChTwoTcErr\,\textdegree C at 10\textdegree C/Ma, considering only the intercept uncertainty. That is the closure temperature of a hematite crystal of about 0.2\,\textmu m radius, or an apatite crystal of about 650\,\textmu m radius.

Closure temperature is the wrong measure for a mineral that forms at the surface and stays there, since it presumes a cooling rate \citep{dodson1973}. Isothermal retention is more appropriate \citep{wolf1998}. For a domain with the parameters of the Stage 1 regression held isothermally at 40\textdegree C for 100\,Ma, $t'/t = 0.968$, where $t'$ is the helium age, $t$ the holding time, and $t'/t = 1$ is perfect retention. At 25\textdegree C the same domain is nearly quantitatively retentive at billion-year timescales.

Comparing whole experiments across samples is easier if the heating schedule is collapsed to a dimensionless diffusive timescale, the Fourier number $Fo = Dt/a^{2}$, calculated against a common reference line. For a single spherical domain the released fraction is a function of Fourier number alone \citep{crank1979,fechtigkalbitzer1966}, so if the reference activation energy matches that of the controlling process, curves for different samples retain their relative positions regardless of the intercept chosen. In this space, vacuum step heats of goethite and of hematite form well separated populations, with goethite the less retentive, and HT-Win-06-01b is the most retentive polycrystalline sample, at least over the 14.7\% of its \textsuperscript{3}He held in the least retentive domains. This suggests that the helium retentivity of goethite that is not decomposing may exceed that of hematite of comparable size, though experiments on irradiated single-crystal goethite of controlled size would be needed to test it. Vacuum step heats of unrelated goethite samples resemble one another more closely than step heats of the same sample at different water activities.

\subsection{Radiogenic helium loss recorded by \textsuperscript{4}He/\textsuperscript{3}He ratios}\label{sec:discussion-ratios}

Bulk-normalized \textsuperscript{4}He/\textsuperscript{3}He ratios in step heats of irradiated samples record the rounding of the radiogenic concentration profile by natural loss, provided loss in nature and in the laboratory propagate with the same geometry \citep{shuster2004,farley2024}. The steep initial rise and subsequent flattening seen in all three Win-06-01b experiments (Fig.~\ref{fig:ratio-spectra}) is consistent with complete, then partial, natural loss from the least retentive domains.

The hydrothermal experiment records lower ratios at any given cumulative \textsuperscript{3}He yield and approaches its plateau most slowly. The three experiments differ systematically in water pressure: \citet{farley2024} heated open to a liquid nitrogen trap, \citet{monteiro2017} at higher water pressures, and this experiment at water saturation. The ordering of the ratio spectra follows that ordering, which suggests the spectra are shaped in part by the extent and spatial propagation of dehydration rather than by natural helium loss alone. The hydrothermal experiment did not reach a plateau and might have with further heating.

Both vacuum experiments show an elevated first-step ratio that is absent here. It represents a negligible fraction of the total \textsuperscript{4}He and may reflect physisorbed atmospheric helium \citep{protin2016}, which would have been removed during the four days at 70\textdegree C preceding the hydrothermal experiment.

% =============================================================
\section{Conclusions}\label{sec:conclusions}
% =============================================================

Sealing a sample in a pure silica glass ampoule preserves reactive phases through helium step heating and allows the experiment to be run under controlled water or sulfur fugacity in a standard ultra-high-vacuum heating cell, at internal water pressures up to about 86\,bar at 300\textdegree C. Helium permeates the ampoule wall as a first-order process, and the lag this introduces between release and measurement can be removed either by measuring the ampoule's residual inventory with low-temperature permeation steps or, for long steps, by bounding it. Encapsulated Durango apatite recovers its published Arrhenius array under this correction. The ampoule's permeation rate constant can be measured during the experiment itself, which is necessary for hydrothermal work because water-bearing ampoules do not follow published values for anhydrous glass.

The first hydrothermal helium step heat of goethite gives apparent diffusivities approximately seven natural log units below those obtained from the same sample in vacuum. The corrected \textsuperscript{3}He array is linear above 7\% cumulative yield, with an activation energy of \ChTwoEa\ $\pm$ \ChTwoEaErr\,kJ/mol, and its delta plot follows multiple diffusion domain systematics. Extrapolated to surface conditions, a domain with these parameters is nearly quantitatively retentive of helium at billion-year timescales at 25\textdegree C.

The immediate consequence is for how vacuum step-heat data on goethite should be read. The difference between the two experiments on the same aliquot of the same sample is far larger than the differences between vacuum experiments on unrelated samples, which indicates that vacuum step heats of goethite are dominated by the response of the mineral to the experiment rather than by helium migration through it. Helium retention in natural goethite is correspondingly more robust than vacuum experiments imply.

We cannot exclude a contribution from dehydration or recrystallization under hydrothermal conditions. If either contributed, the diffusivities reported here are upper bounds and the retentivity a lower bound, which would strengthen rather than weaken the conclusion. Distinguishing these possibilities, and separating the effects of crystallite size, chemical composition, and radiation damage, will require hydrothermal experiments on irradiated single-crystal goethite of controlled size. The method described here makes those experiments straightforward.
